
\chapter{Ring Theory}

\section{Ideal}

\section{Ring of Fractions}
\begin{theorem}
    Let $R$ be a Commutative Ring, $D \subset R$ be a subset such that $\begin{cases}
        \text{no zero, no zero divisors} \\ \text{closed under multiplication}
    \end{cases}$. \\
    Then, there exists a Commutative Ring $Q$ with identity satisfies:
    \begin{enumerate}
        \item $R$ can embed in $Q$, and every element of $D$ becomes unit in $Q$.
        More precisely, $Q = \{r d^{-1} \mid r \in R,\ d \in D\}$.
        \item $Q$ is the smallest Ring contaning $R$ with identity such that every element of $D$ becomes unit in $Q$.
    \end{enumerate}
\end{theorem}
\begin{proof}
    Let $\mathcal{F} \defequal \{(r,d) \mid r \in R,\ d \in D\}$ and the relation $\sim$ on $\mathcal{F}$
    by $(r_1, d_1) \sim (r_2, d_2) \iff r_1d_2 = r_2d_1$.\\
    Then, $\sim$ is equivalent relation: reflexive and symmetirc are clear, and
    Suppose that $(r_1, d_1) \sim (r_2, d_2)$ and $(r_2, d_2) \sim (r_3, d_3)$.
    \begin{align*}
        r_2d_3 = r_3d_2 \implies r_2 d_1 d_3 = r_3 d_1 d_2 \implies r_1 d_2 d_3 = r_3 d_1 d_2 \implies d_2 ( r_1 d_3 - r_3 d_1) \implies r_1d_3 = r_3d_1
    \end{align*}
    Thus transitivity shown. Define
    \begin{align*}
        \frac{r}{d} \defequal [(r,d)] = \{(a,b) \mid (a,b) \sim (r,d) \},\ Q \defequal \left\{ \frac{r}{d} \;\middle|\; r \in R,\ d \in D \right\}
    \end{align*}
    And define operations $+, \times$ on $Q$:
    \begin{align*}
        \frac{r_1}{d_1} + \frac{r_2}{d_2} \defequal \frac{r_1d_2 + r_2d_1}{d_1 d_2}, \quad
        \frac{r_1}{d_1} \times \frac{r_2}{d_2} \defequal \frac{r_1r_2}{d_1d_2}
    \end{align*}
    Well-Definedness: If $\dis \frac{r_1}{d_1} = \frac{r'_1}{d'_1}$ and $\dis \frac{r_2}{d_2} = \frac{r'_2}{d'_2}$,
    \begin{align*}
        \frac{r_1d_2 + r_2d_1}{d_1 d_2}
        = \frac{r_1 d_2 d'_1 d'_2 + r_2 d_1 d'_1 d'_2}{d_1 d_2 d'_1 d'_2}
        = \frac{(r_1 d'_1) d_2 d'_2 + (r_2 d'_2) d_1 d'_1 }{d_1 d_2 d'_1 d'_2}
        = \frac{(r'_1 d_1) d_2 d'_2 + (r'_2 d_2) d_1 d'_1 }{d_1 d_2 d'_1 d'_2}
        = \frac{(r'_1d'_2 + r'_2d'_1)d_1d_2}{d_1d_2 d'_1 d'_2}
        = \frac{r'_1d'_2 + r'_2d'_1}{d'_1 d'_2}
    \end{align*}
    \begin{align*}
        \frac{r_1r_2}{d_1d_2} 
        = \frac{r_1 r_2 d'_1 d'_2}{d_1d_2d'_1d'_2}
        = \frac{(r_1 d'_1) (r_2 d'_2)}{d_1d_2d'_1d'_2}
        = \frac{(r'_1 d_1) (r'_2 d_2)}{d_1d_2d'_1d'_2}
        = \frac{r'_1 r'_2 d_1 d_2}{d_1d_2d'_1d'_2}
        = \frac{r'_1 r'_2}{d'_1d'_2}
    \end{align*}
    Now, $(Q, +, \times)$ constructs Commutative Ring with identity: for any $d \in D$, 
    put $\dis 0_Q \defequal \frac{0}{d},\ 1_Q \defequal \frac{d}{d}$. Then,
    \begin{enumerate}
        \item $(R, +, \times)$ closed under the operations since $D$ is closed under the multiplication.
        \item $(R, +)$ has a zero: $\dis \frac{r_1}{d_1} + 0_Q = \frac{r_1}{d_1} + \frac{0}{d} = \frac{r_1 d + 0 d_1}{d_1 d} = \frac{r_1 d}{d_1 d} = \frac{r_1}{d_1}$.
        \item $(R, +)$ has an inverse: $\dis \frac{r_1}{d_1} + \frac{-r_1}{d_1} = \frac{r_1d_1 + (-r_1)d_1}{d_1 d_1} = \frac{[(r_1) + (-r_1)]d_1}{d_1 d_1} = \frac{0 d_1}{d_1 d_1} = \frac{0}{d_1 d_1} = 0_Q$.
        \item $(R, +, \times)$ satisfies distributive law:
        \begin{enumerate}[label=\theenumi-\arabic*.]
            \item The left law:
            \begin{align*}
                \frac{r_1}{d_1} \times \left( \frac{r_2}{d_2} + \frac{r_3}{d_3} \right)
            = &\frac{r_1}{d_1} \times \frac{r_2 d_3 + r_3 d_2}{d_2 d_3}
            = \frac{r_1 r_2 d_3 + r_1 r_3 d_2}{d_1 d_2 d_3}
            = \frac{r_1 r_2 d_1 d_3 + r_1 r_3 d_1 d_2}{d_1 d_2 d_1 d_3}
            = \frac{r_1 r_2}{d_1 d_2} + \frac{r_2 r_3}{d_2 d_3} \\
            = &\frac{r_1}{d_1} \times \frac{r_2}{d_2} + \frac{r_2}{d_2} \times \frac{r_3}{d_3}
            \end{align*}
            \item The right law:
            \begin{align*}
                \left(\frac{r_1}{d_1} + \frac{r_2}{d_2}\right) \times \frac{r_3}{d_3}
            = &\frac{r_1 d_2 + r_2 d_1}{d_1 d_2} \times \frac{r_3}{d_3}
            = \frac{r_1 r_3 d_2 + r_2 r_3 d_1}{d_1 d_2 d_3}
            = \frac{r_1 r_3 d_2 d_3 + r_2 r_3 d_1 d_3}{d_1 d_3 d_2 d_3}
            = \frac{r_1 r_3}{d_1 d_3} + \frac{r_2 r_3}{d_2 d_3} \\
            = &\frac{r_1}{d_1} \times \frac{r_3}{d_3} + \frac{r_2}{d_2} \times \frac{r_3}{d_3}
            \end{align*}
        \end{enumerate}
        \item $(R, \times)$ has an identity: $\dis \frac{r_1}{d_1} \times 1_Q = \frac{r_1}{d_1} \times \frac{d}{d} = \frac{r_1d}{d_1d} = \frac{r_1}{d_1}$.
        \item Elements of $D$ become unit in $Q$: Define $\iota:R \to Q:r \mapsto \dis \frac{rp}{p}$ where $p \in D$ is any fixed element in $D$. \\
        Then, $\iota$ is Ring-Monomorphsim because:
        \begin{enumerate}[label=\theenumi-\arabic*.]
            \item Well-Defined and Injective: 
            $\dis \iota (r_1) = \iota (r_2)
            \iff \frac{r_1 p}{p} = \frac{r_2 f}{f}
            \iff (r_1 - r_2) pp = 0
            \iff r_1 = r_2$
            \item For any $d \in D$, $\iota(d)$ is a unit of $Q$: Put $\dis (\iota(d))^{-1} \defequal \frac{p}{dp}$, then
            \begin{align*}
                \iota(d) \times (\iota(d))^{-1} = \frac{dp}{p} \times \frac{p}{dp} = \frac{dpp}{dpp} = 1_Q
            \end{align*}
            \noindent That is, $\iota$ is embedding from $R$ into $Q$ such that $\iota[D]$ becomes units of $Q$ except zero. \\
            Moreover, if $D = R \setminus \{0\}$, then $Q$ is field.
        \end{enumerate}
        \item $Q$ is the \textit{smallest} ring containing $R$ with identity such that every element of $D$ becomes units in $Q$. \\
        Let $S$ be an any commutative ring with identity, \\
        and assume that $\varphi:R \to S$ is a Ring-Monomorphism such that for any $d \in D$, $\varphi(d)$ is unit in $S$. \\
        Define $\dis \phi:Q \to S:\frac{r}{d} \mapsto \varphi(r) \varphi(d)^{-1}$. Then, this $\phi$ is well-defined and injective:
        \begin{align*}
            &\phi \left(\frac{r_1}{d_1} \right) = \phi \left(\frac{r_2}{d_2} \right)
            \iff
            \varphi(r_1) \varphi(d_1)^{-1} = \varphi(r_2) \varphi(d_2)^{-1}
            \iff 
            \varphi(r_1) \varphi(d_2) = \varphi(r_2) \varphi(d_1)\\
            \overset{\text{homomor.}}{\iff}
            &\varphi(r_1 d_2) = \varphi(r_2 d_1) 
            \overset{\text{one-to-one}}{\iff}
            r_1 d_2 = r_2 d_1
            \iff
            \frac{r_1}{d_1} = \frac{r_2}{d_2}
        \end{align*}
        That is, if a commutative ring $S$ with identity contains a copy of $R$ such that the denominator set $D$ of $R$ becomes unit in $S$, then $S$ contains ring of fractions $Q$ of $R$.
        Thus $S = Q$ is the smallest ring that satisfies these conditions.
    \end{enumerate}
\end{proof}

\section{Commutative Ring with Unity}

\chapter{Polynomial Ring Theory}